\documentclass[10pt,a4paper]{amsart}
\usepackage{amsmath,amssymb,amsthm,graphicx}
\usepackage{verbatim}
\newenvironment{reference}{\comment}{\endcomment}
%\newenvironment{reference}{}{}
\newenvironment{slogan}{\comment}{\endcomment}
\newenvironment{history}{\comment}{\endcomment}
% For commutative diagrams we use Xy-pic
\usepackage[all]{xy}

% We use 2cell for 2-commutative diagrams.
\xyoption{2cell}
% Theorem environments.
%
\theoremstyle{plain}
\newtheorem{theorem}[subsection]{Theorem}
\newtheorem{proposition}[subsection]{Proposition}
\newtheorem{lemma}[subsection]{Lemma}

\theoremstyle{definition}
\newtheorem{definition}[subsection]{Definition}
\newtheorem{example}[subsection]{Example}
\newtheorem{exercise}[subsection]{Exercise}
\newtheorem{situation}[subsection]{Situation}

\theoremstyle{remark}
\newtheorem{remark}[subsection]{Remark}
\newtheorem{remarks}[subsection]{Remarks}

\numberwithin{equation}{subsection}

% Macros
%
\def\lim{\mathop{\mathrm{lim}}\nolimits}
\def\colim{\mathop{\mathrm{colim}}\nolimits}
\def\Spec{\mathop{\mathrm{Spec}}}
\def\Hom{\mathop{\mathrm{Hom}}\nolimits}
\def\Ext{\mathop{\mathrm{Ext}}\nolimits}
\def\SheafHom{\mathop{\mathcal{H}\!\mathit{om}}\nolimits}
\def\SheafExt{\mathop{\mathcal{E}\!\mathit{xt}}\nolimits}
\def\Sch{\mathit{Sch}}
\def\Mor{\mathop{\mathrm{Mor}}\nolimits}
\def\Ob{\mathop{\mathrm{Ob}}\nolimits}
\def\Sh{\mathop{\mathit{Sh}}\nolimits}
\def\NL{\mathop{N\!L}\nolimits}
\def\CH{\mathop{\mathrm{CH}}\nolimits}
\def\proetale{{pro\text{-}\acute{e}tale}}
\def\etale{{\acute{e}tale}}
\def\QCoh{\mathit{QCoh}}
\def\Ker{\mathop{\mathrm{Ker}}}
\def\Im{\mathop{\mathrm{Im}}}
\def\Coker{\mathop{\mathrm{Coker}}}
\def\Coim{\mathop{\mathrm{Coim}}}

% Boxtimes
%
\DeclareMathSymbol{\boxtimes}{\mathbin}{AMSa}{"02}

%
% Macros for moduli stacks/spaces
%
\def\QCohstack{\mathcal{QC}\!\mathit{oh}}
\def\Cohstack{\mathcal{C}\!\mathit{oh}}
\def\Spacesstack{\mathcal{S}\!\mathit{paces}}
\def\Quotfunctor{\mathrm{Quot}}
\def\Hilbfunctor{\mathrm{Hilb}}
\def\Curvesstack{\mathcal{C}\!\mathit{urves}}
\def\Polarizedstack{\mathcal{P}\!\mathit{olarized}}
\def\Complexesstack{\mathcal{C}\!\mathit{omplexes}}
% \Pic is the operator that assigns to X its picard group, usage \Pic(X)
% \Picardstack_{X/B} denotes the Picard stack of X over B
% \Picardfunctor_{X/B} denotes the Picard functor of X over B
\def\Pic{\mathop{\mathrm{Pic}}\nolimits}
\def\Picardstack{\mathcal{P}\!\mathit{ic}}
\def\Picardfunctor{\mathrm{Pic}}
\def\Deformationcategory{\mathcal{D}\!\mathit{ef}}
\usepackage{cleveref}
\crefname{originallemma}{original lemma}{original lemmas}
\Crefname{originallemma}{Original lemma}{Original lemmas}
\makeatletter
\newlabel{lemma-closed-in-quasi-compact}{{lemma-closed-in-quasi-compact}{source}{}{sourcepointer-1}{}}
\newlabel{lemma-closed-in-quasi-compact@cref}{{[originallemma][0][]lemma-closed-in-quasi-compact}{[1][0][]source}{}{}{}}
\makeatother
\crefname{originallemma}{original lemma}{original lemmas}
\Crefname{originallemma}{Original lemma}{Original lemmas}
\makeatletter
\newlabel{lemma-constructible-hausdorff-quasi-compact}{{lemma-constructible-hausdorff-quasi-compact}{source}{}{sourcepointer-2}{}}
\newlabel{lemma-constructible-hausdorff-quasi-compact@cref}{{[originallemma][0][]lemma-constructible-hausdorff-quasi-compact}{[1][0][]source}{}{}{}}
\makeatother
\crefname{originallemma}{original lemma}{original lemmas}
\Crefname{originallemma}{Original lemma}{Original lemmas}
\makeatletter
\newlabel{lemma-image-quasi-compact}{{lemma-image-quasi-compact}{source}{}{sourcepointer-3}{}}
\newlabel{lemma-image-quasi-compact@cref}{{[originallemma][0][]lemma-image-quasi-compact}{[1][0][]source}{}{}{}}
\makeatother
\crefname{originallemma}{original lemma}{original lemmas}
\Crefname{originallemma}{Original lemma}{Original lemmas}
\makeatletter
\newlabel{lemma-quasi-compact-in-Hausdorff}{{lemma-quasi-compact-in-Hausdorff}{source}{}{sourcepointer-4}{}}
\newlabel{lemma-quasi-compact-in-Hausdorff@cref}{{[originallemma][0][]lemma-quasi-compact-in-Hausdorff}{[1][0][]source}{}{}{}}
\makeatother
\title{A spectral map is constructibly continuous, has quasi-compact fibres and constructibly closed image}
\author{The Stacks Project Authors}
\date{stacks-project:topology.tex:local-snapshot-20261012; GNU FDL 1.2 or later}
\begin{document}
\maketitle
\noindent\textit{Editorial note.} All delivered mathematical text is the author's own, copied from the pinned Stacks Project chapter source sources/stacks-topology/pinned-original-source.tex; no statement or proof step was rewritten, reordered or removed.\par\medskip
\noindent\textit{Editorial note.} Preamble provenance: the Stacks Project keeps one single shared preamble for every chapter; this card takes it directly from the staged sources/stacks-topology/preamble.tex as three verbatim spans (lines 8-12, 14-18 and 154-217), so no derived file was written and nothing was added or redefined. The wrapper class and the base packages are supplied by the assembler.\par\medskip
\noindent\textit{Editorial note.} Bibliography provenance: neither the counted statement nor its proof contains a citation (this lane rescans the whole span before writing the card), so the card carries no bibliography and no entry was reconstructed.\par\medskip
\noindent\textit{Editorial note.} Selection provenance: this card was opened under the file-based control contract (CONTROL.md), with the mechanical pre-card gate run on its statement and proof interval before the card was written.\par\medskip
\noindent\textit{Allowed standard prerequisites (editorial).}\begin{itemize}
\item The statement that the constructible topology of a spectral space is Hausdorff and quasi-compact. Invoked at the start of the counted proof (source lines 4599-4601). Puts both spaces in the setting of quasi-compact Hausdorff spaces.
\item The statement that a quasi-compact subset of a Hausdorff space is closed. Invoked inside the counted proof (source lines 4605-4606). Gives closedness of the constructible image.
\item The statement that a closed subset of a quasi-compact space is quasi-compact. Invoked inside the counted proof (source lines 4613-4614). Makes the constructible fibres quasi-compact.
\item The statement that quasi-compactness transfers across a surjective continuous map. Invoked at the end of the counted proof (source line 4616). Gives quasi-compactness of the fibres in the given topology.
\end{itemize}\medskip
\ifdefined\hypertarget\hypertarget{sourcepointer-1}{}\fi
\noindent\textit{Source reference.} Original lemma lemma-closed-in-quasi-compact: sources/stacks-topology/pinned-original-source.tex, source lines 1879--1879. Navigation pointer only: the counted text refers to the Stacks item lemma-closed-in-quasi-compact, which is an existing labeled result of this same chapter and is not part of this card. The pointer carries the target's own label as its number and its definition line as its span; its content is not reproduced here.\par\medskip
\ifdefined\hypertarget\hypertarget{sourcepointer-2}{}\fi
\noindent\textit{Source reference.} Original lemma lemma-constructible-hausdorff-quasi-compact: sources/stacks-topology/pinned-original-source.tex, source lines 4522--4522. Navigation pointer only: the counted text refers to the Stacks item lemma-constructible-hausdorff-quasi-compact, which is an existing labeled result of this same chapter and is not part of this card. The pointer carries the target's own label as its number and its definition line as its span; its content is not reproduced here.\par\medskip
\ifdefined\hypertarget\hypertarget{sourcepointer-3}{}\fi
\noindent\textit{Source reference.} Original lemma lemma-image-quasi-compact: sources/stacks-topology/pinned-original-source.tex, source lines 1964--1964. Navigation pointer only: the counted text refers to the Stacks item lemma-image-quasi-compact, which is an existing labeled result of this same chapter and is not part of this card. The pointer carries the target's own label as its number and its definition line as its span; its content is not reproduced here.\par\medskip
\ifdefined\hypertarget\hypertarget{sourcepointer-4}{}\fi
\noindent\textit{Source reference.} Original lemma lemma-quasi-compact-in-Hausdorff: sources/stacks-topology/pinned-original-source.tex, source lines 1894--1894. Navigation pointer only: the counted text refers to the Stacks item lemma-quasi-compact-in-Hausdorff, which is an existing labeled result of this same chapter and is not part of this card. The pointer carries the target's own label as its number and its definition line as its span; its content is not reproduced here.\par\medskip
\begin{lemma}
\label{lemma-fibres-spectral-map-quasi-compact}
Let $f : X \to Y$ be a spectral map of spectral spaces. Then
\begin{enumerate}
\item $f$ is continuous in the constructible topology,
\item the fibres of $f$ are quasi-compact, and
\item the image is closed in the constructible topology.
\end{enumerate}
\end{lemma}

\begin{proof}
Let $X'$ and $Y'$ denote $X$ and $Y$ endowed with the constructible topology
which are quasi-compact Hausdorff spaces by
Lemma \ref{lemma-constructible-hausdorff-quasi-compact}.
Part (1) says $X' \to Y'$ is continuous and follows immediately from the
definitions. Part (3) follows as $f(X')$ is a quasi-compact subset of the
Hausdorff space $Y'$, see Lemma \ref{lemma-quasi-compact-in-Hausdorff}.
We have a commutative diagram
$$
\xymatrix{
X' \ar[r] \ar[d] & X \ar[d] \\
Y' \ar[r] & Y
}
$$
of continuous maps of topological spaces. Since $Y'$ is Hausdorff
we see that the fibres $X'_y$ are closed in $X'$. As $X'$ is quasi-compact
we see that $X'_y$ is quasi-compact
(Lemma \ref{lemma-closed-in-quasi-compact}).
As $X'_y \to X_y$ is a surjective continuous map we conclude
that $X_y$ is quasi-compact (Lemma \ref{lemma-image-quasi-compact}).
\end{proof}

\end{document}
